\subsubsection{Image and Video}
\label{subsubsection::imagevideo}

Image and video data curation follows a set of carefully designed processing, annotation, and filtering pipelines: (1) collecting raw data and performing pre-processing; (2) computing embeddings and conducting deduplication; (3) categorizing samples and applying basic filtering; (4) annotating data; and (5) grouping samples into training-ready shards based on their resolution and duration. To improve generation quality for Physical AI scenarios and other challenging cases, we introduce synthetic data into the visual-generation data mixture. We organize the resulting data into pre-training data and higher-quality mid-training and post-training data.

\paragraph{Pre-training.}
In the pre-training stage, we use $767$M images and $347.7$M video clips processed from $7.8$B raw images and $3$B raw source videos. In the resulting corpus, $720$p and $480$p are the dominant resolutions for both images and videos. Specifically, $720$p accounts for $26.8\%$ of images and $36.4\%$ of videos, while $480$p accounts for $26.0\%$ of images and $30.8\%$ of videos. In addition, $25.2\%$ of images and $12.2\%$ of videos are at 1080p resolution or higher. The most common aspect ratio is 16:9, accounting for $52.0\%$ of images and $97.3\%$ of videos. For images, the second most common aspect ratio is 1:1, accounting for $25.2\%$ of the retained image corpus. The raw data is processed and filtered using the pipeline described below:

\begin{itemize}

\item \textit{Raw data collection and processing.} We collect billions of raw images and videos from diverse data sources, recording the raw media content together with associated metadata such as raw captions and descriptions. For videos, we additionally apply scene-change detection using TransNetV2~\citep{soucek2020transnetv2} to segment long videos into temporally consistent clips. We then use \texttt{ffmpeg cropdetect} to detect and remove black borders, and re-encode all video clips into a canonical format to standardize storage and ensure playback integrity.

\item \textit{Embedding and deduplication.} Raw data contains a large amount of repeated image and video content, and its concept distribution is often highly imbalanced. To remove duplicate content and establish a foundation for concept balancing and evaluation, we embed the media content into vectors. For images, we use Qwen3-VL-Embedding-8B~\citep{qwen3vlembedding}; for videos, we use nvidia/Cosmos-Embed1-448p~\citep{nvidia2025cosmosembed1}. We sample $147$M images and $400$M video clips from the full data corpus and independently run cuML KMeans with $20,000$ clusters for each data type \citep{cuml, kmeans}. We then assign each image or video clip to its nearest cluster and perform near-duplicate removal within each cluster based on cosine similarity scores.

\item \textit{Categorization and basic filtering.} We use a small suite of in-house VLM models for semantic tagging and quality filtering. Both image and video data are classified into $47$ hierarchical categories, including General and Physical AI domains. For image filtering, a dedicated model annotates attributes such as collage and produces aesthetic and photorealism scores. Images are retained only if their aesthetic score exceeds a predefined threshold. Images tagged as collage, watermark, white background, or NSFW are discarded. For synthetic images not intended for text rendering, we additionally filter them based on their photorealism score, retaining only those above a specified threshold. For video filtering, we use three continuous quality scores---DOVER aesthetic quality~\citep{dover}, DOVER technical quality, and VTSS training suitability~\citep{koala}, each on a 0--9 scale---together with approximately $100$ binary artifact tags. Major artifacts (split-screen layouts, rotated videos, static videos) lead to rejection, while minor artifacts (text overlays, motion blur, compression noise) are flagged but retained in the pre-training data.

\end{itemize}

\paragraph{Mid-training.} The mid-training stage aims to improve generation quality using carefully selected high-quality data and to equip the model with additional capabilities, including both domain-specific capabilities, such as Physical AI, and new tasks, such as video transfer. The data is drawn from three sources: (1) high-quality images and videos; (2) synthetic images and videos; and (3) video-transfer data.

\begin{itemize}
\item \textit{High-quality images and videos.} We select data using stricter filtering rules and sample data more heavily from hard-case concepts to mitigate the long-tailed distribution. For real images, samples must satisfy per-aspect-ratio resolution thresholds and a strict DOVER aesthetic-score cutoff. We also include synthetic and text-rendering subsets. Synthetic images, curated through careful rejection sampling, broaden coverage of uncommon visual concepts and object compositions, while text-rendering images address the underrepresentation of legible in-image text. The resulting mid-training image mixture has effective proportions of $60\%$ real images, $36\%$ synthetic images, and $4\%$ text-rendering images. For video, we similarly apply stricter resolution and aesthetic filtering to select clips from the pre-training pool, which constitute $46.0\%$ of the mid-training video mixture. We then incorporate additional high-quality, domain-specific clips from robotics, autonomous driving, human activity, and egocentric human-object interaction sequences, targeting embodied-AI and manipulation scenarios; these clips constitute another $43.9\%$ of the mixture. To further improve robustness on difficult and corner-case concepts, such as human motion, high-speed complex motion, and fine-grained manipulation, we collect additional capability-oriented data focused on these hard cases, which accounts for the remaining $10.1\%$ of mid-training videos.

\item \textit{Synthetic data.} Although our pre-training corpus is highly diverse, its concept distribution remains long-tailed. As a result, the model receives comparatively limited exposure to rare but important Physical AI domains and scenarios, such as robotics, autonomous driving, and warehouse environments. In these settings, the model often struggles to understand scene dynamics, physical interactions, and long-horizon behavior. To address these limitations, we construct a large-scale synthetic data corpus with the following subsets: 1) Physical-Interaction-Scenes \textit{(SDG-PhyxSim)} focusing on rigid-body collisions, articulated object dynamics, deformable materials, fluid dynamics, and optical effects; 2) Embodied-Robot-Scenes \textit{(SDG-RobotSim)} for manipulation and locomotion sequences across 6--8 robot embodiments and diverse task categories; 3) Autonomous-Driving-Scenarios \textit{(SDG-DriveSim)} covering both routine and corner-case traffic scenarios; 4) Digital-Human-Scenes \textit{(SDG-SynHuman)} designed to improve modeling of human dynamics, camera-motion priors, and multi-character interactions; and 5) Warehouse-Operation-Scenes \textit{(SDG-Warehouse)} for warehouse safety containing human-forklift interaction scenarios. Refer to Appendix~\ref{appendix:sdg_datasets} for more details on the construction of the synthetic data and analysis. We release all SDG datasets to support the community.

\item \textit{Video transfer data.} Transfer data equips the Generator with control-conditioned generation capabilities. Given a spatial control signal, such as an edge map, blurred frame, depth map, segmentation map, or world-scenario map, together with a text description, the model is trained to generate an RGB video. We select $3$M videos from the pre-training video pools, focusing on high-quality videos and physical-AI domains such as robotics and autonomous driving. For edge and blur control, we compute the control signals on the fly during training using Canny edge detection, Gaussian blur, and bilateral filtering with randomly sampled parameters. For depth and segmentation control, we pre-compute the control signals using Video Depth Anything~\citep{chen2025videodepth} and SAMv2~\citep{ravi2024sam}, respectively. For world-scenario-map control, we use the MADS dataset collected by Cosmos-Drive-Dreams~\citep{ren2025cosmos}. MADS contains $1.1$M samples, each with seven synchronized camera views: front-wide, front-tele, cross-left, cross-right, rear-left, rear-right, and rear-tele. The videos are recorded at $30$\,FPS and accompanied by per-camera world-scenario-map control inputs that encode lane lines, road boundaries, traffic lights, and dynamic 3D bounding boxes for vehicles and pedestrians. The dataset covers 14 geographic regions, including the United States, Germany, Japan, and the United Kingdom, across 25 country-duration partitions.
\end{itemize}

\paragraph{Post-training.} We construct post-training datasets for training domain-specialized Cosmos 3 models, such as Cosmos3-Super-Text2Image and Cosmos3-Super-Image2Video.

The post-training image corpus is a compact, carefully curated set drawn from three sources: synthetic images, text-rendering images, and high-quality real images. General web-scale pre-training data is excluded entirely in favor of high-fidelity content that directly targets generation quality and capability gaps.


The post-training video corpus is assembled from two components. The first is a subsampled subset of the pre-training video corpus.
This subset serve as a regularizer, preventing the model from overfitting to the narrower post-training distribution while preserving the broader visual knowledge acquired during pre-training. The second and primary component is a compact set of supervised fine-tuning (SFT) videos, curated specifically to close generation-quality gaps identified through systematic evaluation. The SFT video set consists of three sub-sources. The first is synthetic videos, which cover diverse visual concepts, motion types, and scene compositions that are difficult to source from real-world footage. The second is human-curated real SFT videos, selected and annotated by human curators to provide a direct quality signal for generation fidelity across key visual domains. The third is retrieved real videos from the pre-training corpus, selected via embedding similarity to ensure coverage of common failure cases.

\paragraph{Structured caption annotation.}
Caption quality is a critical factor in generation quality. A high-quality caption should faithfully describe the entities, attributes, relationships, and overall scene content in an image or video. For videos, captions should further capture temporal dynamics, including object motion, human actions, physical changes, interactions, and camera movement. To improve caption quality, we conducted multiple design iterations and adopted a structured JSON annotation format instead of dense free-form natural-language captions for all our data across all training stages. Our experiments show that free-form captions are often precise but incomplete: they tend to describe visible content accurately, yet omit important details in complex scenes. In contrast, a rich predefined structure encourages systematic coverage of objects, attributes, relationships, and scene-level information, improving recall while maintaining high precision.

Our structured format captures a broad set of visual attributes, including subjects, background, lighting, aesthetics, and cinematography. For videos, we additionally introduce fields for temporal dynamics, ranging from physical transformations and object interactions to complex human motion. We fine-tuned two \mbox{Qwen3-VL-8B} models on structured annotation data to serve as our in-house captioners for images and videos, respectively. Refer to Appendix~\ref{appendix:captioning_details} for more details on our captioning models and full structured caption schema.

To quantitatively and rigorously evaluate annotation quality, we designed a specialized caption-quality benchmark for both images and videos. This benchmark focuses on hard-to-caption examples, emphasizing domains such as Physical AI, where accurate descriptions of objects, spatial relationships, actions, and temporal dynamics are critical. For each model-generated caption, we compute precision and recall at the assertion level using two distinct approaches. Precision is evaluated directly against the source media to penalize hallucinations: a VLM decomposes the generated caption into atomic claims and verifies whether each claim is visually supported by the image or video itself. Recall, on the other hand, measures comprehensiveness and relies on human-curated ground truth. To enable a reliable and traceable recall evaluation, we decompose the visual content of videos or images into a list of atomic assertions covering entities, attributes, relationships, events, and other relevant details. An LLM then cross-references the generated caption against this ground-truth assertion list to determine which key details were successfully captured. This protocol allows us to evaluate not only the factual accuracy of the captions, but also whether they contain the critical visual information required to train high-quality generative models. On this benchmark, our structured annotation approach significantly improved recall while maintaining high precision.
